% -----------------------------------------------------------------------------
% This file is automatically generated.
% Do not edit manually.
% Source: notebooks/support/regression-analysis/support.Rmd
% -----------------------------------------------------------------------------


\noindent This exercise is also an illustration rather than
a case step. It stays in the student notebook
because it introduces the visual distinction
between a confidence interval for the mean
response and a prediction interval for a new
observation before we return to the case-based
forecasting exercises.
Figure 5.18 shows a regression line with a dotted
confidence interval and a shaded prediction interval. It is
created with the following code.

\par\addvspace{\topsep}
\begingroup
\fvset{listparameters={%
  \setlength{\topsep}{0pt}%
  \setlength{\partopsep}{0pt}%
  \setlength{\parsep}{0pt}%
  \setlength{\itemsep}{0pt}%
}}
\begin{Verbatim}[breaklines=true,formatcom=\color{ada_blue}]
from IPython.display import display
alpha_pred = 0.01
t_value_pred = t_dist.ppf(1 - alpha_pred / 2, df=n - 2)
se_pred = sigma * np.sqrt(1 + 1 / n + ((x_extended - 
mean_x) ** 2) / S_xx)
pi_upper_true = (beta_0 + beta_1 * x_extended) + 
t_value_pred * se_pred
pi_lower_true = (beta_0 + beta_1 * x_extended) - 
t_value_pred * se_pred
pi_bounds_true = pd.DataFrame({'x': x_extended, 'pi_lower':
 pi_lower_true, 'pi_upper': pi_upper_true})
fig, ax = plt.subplots(figsize=(10, 6))
display(ax.fill_between(pi_bounds_true['x'], 
pi_bounds_true['pi_lower'], pi_bounds_true['pi_upper'], 
color='#00338D', alpha=0.15))
display(ax.plot(ci_bounds_true['x'], 
ci_bounds_true['ci_lower'], color='#00338D', linewidth=1, 
linestyle='--'))
display(ax.plot(ci_bounds_true['x'], 
ci_bounds_true['ci_upper'], color='#00338D', linewidth=1, 
linestyle='--'))
display(ax.plot(true_line['x'], true_line['y'], 
color='#00338D', linewidth=1.5))
ax.set_xlabel('x')
ax.set_ylabel('y')
ax.set_xlim(x_extended_min, x_extended_max)
plt.tight_layout()
plt.show()
\end{Verbatim}
\endgroup
\par\addvspace{\topsep}
